% Generated draft: requires mathematical and source review.
\documentclass[11pt]{article}
\usepackage[margin=25mm]{geometry}
\usepackage{fontspec}
\setmainfont{Latin Modern Roman}
\usepackage{amsmath,amssymb,mathtools,unicode-math}
\setmathfont{Latin Modern Math}
\usepackage{enumitem,needspace,xurl,hyperref}
\setlength{\parindent}{0pt}
\setlength{\parskip}{5pt}
\setlength{\emergencystretch}{3em}
\newcommand{\N}{\mathbb{N}}
\newcommand{\Z}{\mathbb{Z}}
\newcommand{\Q}{\mathbb{Q}}
\newcommand{\R}{\mathbb{R}}
\newcommand{\C}{\mathbb{C}}
\newcommand{\F}{\mathbb{F}}
\newcommand{\A}{\mathbb{A}}
\newcommand{\PP}{\mathbb{P}}
\newcommand{\E}{\mathbb{E}}
\newcommand{\eps}{\varepsilon}
\newcommand{\dd}{\,\mathrm{d}}
\begin{document}
\section*{UnsolvedMath ID: 800082}
\textbf{Source problem code:} catalogue entry
\subsection{Definitions and mathematical statement}
The source asks about a representation \(\rho\colon\pi_1(\Sigma_g)\to\mathrm{Diff}_0(S^1)\) whose images of generators are close to \(\mathrm{PSL}_2(\R)\), and about bounding its Godbillon--Vey (GV) class. The source does not specify the topology or distance measuring closeness, the differentiability class, or how the GV class is to be bounded. Consequently, the following is a conditional formulation requiring these choices to be resolved.

Under the conventional interpretation that \(\Sigma_g\) is a closed, connected, oriented surface of genus \(g\), and that \(\mathrm{Diff}_0(S^1)\) is the identity component of the group of smooth circle diffeomorphisms, form the suspension circle bundle
\[
E_\rho=(\widetilde{\Sigma}_g\times S^1)/\pi_1(\Sigma_g),\qquad
\gamma\cdot(x,t)=(\gamma x,\rho(\gamma)(t)).
\]
The slices \(\widetilde{\Sigma}_g\times\{t\}\) induce a codimension-one foliation \(\mathcal F_\rho\) on \(E_\rho\). Its Godbillon--Vey class is
\[
\operatorname{GV}(\mathcal F_\rho)=[\eta\wedge\dd\eta]\in H^3(E_\rho;\R),
\]
where a nowhere-vanishing defining one-form \(\omega\) and a one-form \(\eta\) satisfy \(\dd\omega=\eta\wedge\omega\). Using the base and fiber orientations gives the associated number
\[
\operatorname{gv}(\rho)=\bigl\langle\operatorname{GV}(\mathcal F_\rho),[E_\rho]\bigr\rangle.
\]
This numerical interpretation of the source's phrase ``bound the GV class'' requires confirmation.

Choose an identification \(S^1\cong\R\mathrm P^1\), so that the projective action realizes \(\mathrm{PSL}_2(\R)\) as a subgroup of circle diffeomorphisms. If \(d\) is the intended distance, write
\[
\delta_d(f)=\inf_{A\in\mathrm{PSL}_2(\R)}d(f,A).
\]
For a generating set \(S\subseteq\pi_1(\Sigma_g)\), the proximity hypothesis can then be expressed as
\[
\delta_d\bigl(\rho(s)\bigr)\leq\epsilon\qquad\text{for every }s\in S.
\]
The proposed question is whether, for sufficiently small \(\epsilon>0\), this hypothesis implies an estimate
\[
\left|\operatorname{gv}(\rho)\right|\leq B(g,\epsilon)
\]
with a finite bound depending only on \(g\) and \(\epsilon\), once the geometric conventions are fixed. The admissible genus range, meaning of ``sufficiently small,'' treatment of the generating set, choice of distance, and intended scalar or cohomological target remain unspecified. The source also does not prescribe sharpness or limiting behavior of \(B(g,\epsilon)\).
\subsection{Short English statement}
Question 23 asks whether the Godbillon--Vey invariant of a circle action of a genus-\(g\) surface group can be bounded in terms of \(g\) and \(\epsilon\) when every generator acts within \(\epsilon\) of a projective transformation. A precise statement requires a specified notion of proximity, conventions for the generators and surface, and a meaning for the size of the GV class.
\subsection{Sources}
\begin{itemize}
\item Local UnsolvedMath registry, supplied source record with permanent ID \(800082\), exact source code \texttt{catalogue entry}, title \emph{Foliations and Diffeomorphism Groups --- Question 23}. The supplied record labels the problem \emph{open}; this label is not an independently verified mathematical or literature assessment.
\item Nariman, Yazdi, \emph{Problem List on Foliations and Diffeomorphism Groups} (2025), Question 23. Supplied source URL: \url{https://arxiv.org/abs/2503.18274}. The record attributes the problem to Nariman and Yazdi, with proposed year 2025.
\item Supplied provenance records access on 2026-10-16 and extraction by \texttt{arxiv\_latex}. These are supplied provenance statements; no source visit or independent verification is asserted. Supplied upstream revision: \texttt{d371ee917ac9647fbb06cd6fc529deec6f3dda82}.
\end{itemize}
\end{document}
